% AUDIT DE VERACITE 2026-09-20 -- corrections des citations de rapport de jury.
% Deux attributions fabriquees au « jury 2023 » autour de « verifier que le resultat est un entier » :
% cette formule est absente des 8 rapports 2018-2025 (la chaine « un entier » n'y figure nulle part).
% Elle provient du corrige technique du Rapport_Jury_2023.md ligne 720 (« ... = 3,16 et entier donc
% p = 3 »), qui est une ligne de calcul et non un commentaire de jury.
%   - Bloc 4, point 2 de la liste : la citation entre guillemets est SUPPRIMEE ; la regle de methode
%     (arrondir p a l'entier le plus proche puis recalculer Ns) est conservee.
%   - Bloc 4, encadre jury : la troisieme citation (« Rapport jury 2023 ») est SUPPRIMEE, ainsi que
%     le \\ de liaison devenu orphelin. Aucun remplacement pertinent trouve.
% Les trois autres citations de la fiche ont ete verifiees conformes et laissees intactes :
% Rapport_Jury_2024.md l.479 et l.511 (encadre jury) et l.515 (point 6 de la liste).
\expandafter\ifx\csname CAPETINCLUDEMODE\endcsname\relax
\documentclass[11pt,a4paper]{article}
% preamble.tex — packages et commandes communes a toutes les fiches CAPET SII IE
\input{glyphtounicode}
\pdfgentounicode=1
\usepackage{cmap}
\usepackage[T1]{fontenc}
\usepackage[utf8]{inputenc}
\usepackage[french]{babel}
% babel[french] active a \begin{document} des raccourcis de ponctuation (:;!?)
% qui cassent la syntaxe to[...] de circuitikz (erreur "+ or - expected") ;
% on les desactive, ce qui ne change que l'espacement fin avant ces signes.
\AtBeginDocument{\shorthandoff{:;!?}}
\usepackage[a4paper, top=1.6cm, bottom=1.6cm, left=1.4cm, right=1.4cm]{geometry}
\setlength{\headheight}{14.5pt}
\usepackage{amsmath, amssymb, amsfonts}
\usepackage{siunitx}
\sisetup{locale=FR, output-decimal-marker={,}, inter-unit-product=\ensuremath{{}\cdot{}},
         group-minimum-digits=4}
% Unites non fournies par siunitx, declarees ici une fois pour toutes (fiches + exercices).
% Les gardes \ifdefined evitent tout conflit avec une redeclaration locale restee dans un
% fichier : ne pas retirer les gardes, ne pas redeclarer ces unites dans les fichiers.
\DeclareSIUnit{\jour}{j}
\ifdefined\var\else\DeclareSIUnit{\var}{var}\fi   % puissance reactive
\ifdefined\vad\else\DeclareSIUnit{\vad}{VAD}\fi   % puissance deformante
\ifdefined\tr\else\DeclareSIUnit{\tr}{tr}\fi      % tour (tr/min)
\usepackage{array, booktabs, multirow, tabularx}
\usepackage{xcolor}
\usepackage[most, breakable]{tcolorbox}
\usepackage{multicol}
\usepackage{enumitem}
\usepackage{mathtools}
\usepackage{lmodern}
% Refonte visuelle 2026 (direction "moderne coloré") : Lato pour le texte,
% maths sans-serif assorties via newtxsf (charge apres amsmath/amssymb).
\usepackage[default]{lato}
\usepackage{newtxsf}
\usepackage{microtype}
\usepackage{fancyhdr}
\usepackage{tikz}
\usetikzlibrary{positioning,arrows.meta}
\usepackage{circuitikz}
\usepackage{lastpage}

% ── Couleurs par domaine (palette modernisée 2026) ────────────────────
\definecolor{colElec}{RGB}{29, 78, 216}        % Électrotechnique — bleu vif
\definecolor{colElecLight}{RGB}{219, 234, 254} % fond bleu clair
\definecolor{colAuto}{RGB}{21, 128, 61}        % Automatique — vert
\definecolor{colMeca}{RGB}{194, 65, 12}        % Mécanique — orange
\definecolor{colTherm}{RGB}{185, 28, 28}       % Thermique — rouge
\definecolor{colQT}{RGB}{232, 244, 233}        % (hérité, plus utilisé par questiontype)
\definecolor{colJury}{RGB}{255, 248, 225}      % fond jury (crème)
\definecolor{colJuryFrame}{RGB}{202, 138, 4}   % bordure/bandeau jury (ambre)
\definecolor{colGrayLight}{RGB}{245, 245, 247} % fond gris neutre
\definecolor{colSeparator}{RGB}{180, 180, 180} % gris séparateur
\definecolor{colExo}{RGB}{8, 145, 178}          % Banque d'exercices — cyan (distinct des fiches)
% ── Encre et gris de la refonte ───────────────────────────────────────
\definecolor{colInk}{RGB}{23, 32, 42}          % texte courant (quasi-noir)
\definecolor{colSlate}{RGB}{83, 98, 118}       % gris ardoise (méta, en-têtes)
\definecolor{colRule}{RGB}{203, 213, 225}      % filets clairs
\definecolor{colJuryBg}{RGB}{255, 248, 225}    % fond crème jury
\definecolor{colPresence}{RGB}{83, 98, 118}    % bloc 5 (présence dans les sujets)
\AtBeginDocument{\color{colInk}}

% ── Style encadré exercice + corrigé (banque d'exercices) ─────────────
\tcbset{
  exoBox/.style={
    breakable,
    enhanced,
    colback=colExo!8!white,
    colframe=colExo,
    colbacktitle=colExo,
    coltitle=white,
    fonttitle=\bfseries\small,
    left=4pt, right=4pt, top=3pt, bottom=3pt,
    before skip=5pt, after skip=5pt,
  }
}

% ── Bascule enonce / corrige (banque d'exercices) ─────────────────────
% \CAPETCORRIGE est defini en ligne de commande (pdflatex "\def\CAPETCORRIGE{1}\input{...}")
% pour compiler la variante corriges d'un meme fichier source ; par defaut (non defini),
% la compilation produit la variante enonces. \CAPETSHOW{enonce}{corrige} choisit la branche.
%
% TROIS MODES (refonte 2026-09-20, manuel de revision) :
%   - aucun drapeau   -> imprime l'enonce            (variante enonces, INCHANGE)
%   - \CAPETCORRIGE   -> imprime le corrige          (variante corriges, INCHANGE)
%   - \CAPETMANUEL    -> imprime l'enonce ET met le corrige de cote pour l'annexe
%
% Le mode manuel stocke chaque corrige dans une macro (\csgdef), sans ecriture
% fichier ni re-tokenisation : le contenu (circuitikz, accents, %) traverse intact.
% L'annexe les rejoue ensuite dans l'ordre via \CAPETANNEXE.
\usepackage{etoolbox}

\newcounter{CAPETexo}          % numero continu d'exercice dans le manuel
\setcounter{CAPETexo}{0}
\newcounter{CAPETannexe}       % compteur de relecture en annexe

% Titre courant de la notion en cours (pose par \CAPETNOTION), sert d'intertitre annexe
\newcommand{\CAPETnotioncode}{}
\newcommand{\CAPETnotiontitre}{}

% \CAPETNOTION{code}{titre} : ouvre une notion dans le manuel (memorise le libelle,
% et marque le rang du premier exercice de la notion pour l'intertitre d'annexe).
\newcommand{\CAPETNOTION}[2]{%
  \renewcommand{\CAPETnotioncode}{#1}%
  \renewcommand{\CAPETnotiontitre}{#2}%
  \csgdef{CAPETnotionstart\the\numexpr\value{CAPETexo}+1\relax}{#1 -- #2}%
}

% ── \degree robuste ───────────────────────────────────────────────────
% Certains fichiers d'exercices utilisent \degree. Ni le preambule ni siunitx
% ne le definissent, et \providecommand local ne suffit pas : sous hyperref,
% \degree existe deja comme commande factice (\ERROR), donc \providecommand
% s'abstient et l'appel echoue a la compilation du manuel (erreur invisible
% en compilation isolee). On le (re)definit donc inconditionnellement ici,
% apres le chargement des packages, pour tout le corpus.
\AtBeginDocument{\providecommand{\degree}{}\renewcommand{\degree}{\ensuremath{^\circ}}}

% Libelle decoratif : les fichiers d'exercices utilisent aussi \CAPETSHOW dans
% leur \fancyhead et leur bloc de titre, pour afficher "(enonces)" ou "(corriges)".
% Ces appels-la ne sont PAS des exercices : en mode manuel ils ne doivent ni
% consommer un numero, ni alimenter l'annexe. On les reconnait a leur premier
% argument, strictement egal au libelle ci-dessous.
\newcommand{\CAPETLIBELLEENONCES}{(\'enonc\'es)}

\ifdefined\CAPETMANUEL
  % --- Mode manuel : enonce en place, corrige differe vers l'annexe ---
  \newcommand{\CAPETSHOW}[2]{%
    \def\CAPETargun{#1}%
    \ifx\CAPETargun\CAPETLIBELLEENONCES
      #1%            appel decoratif : rendu tel quel, hors numerotation
    \else
      \CAPETSHOWEXO{#1}{#2}%
    \fi
  }
  \newcommand{\CAPETSHOWEXO}[2]{%
    \stepcounter{CAPETexo}%
    \csgdef{CAPETcorrige\theCAPETexo}{#2}%
    \edef\CAPETtmpnum{\theCAPETexo}%
    #1%
    \par\vspace{2pt}%
    {\footnotesize\itshape\color{colExo}%
     \hyperref[capet:cor:\CAPETtmpnum]{Corrige de l'exercice \CAPETtmpnum\ $\rightarrow$ p.~\pageref{capet:cor:\CAPETtmpnum}}%
     \label{capet:eno:\CAPETtmpnum}}%
  }
  % \CAPETANNEXE : rejoue tous les corriges captures, dans l'ordre, avec intertitres
  \newcommand{\CAPETANNEXE}{%
    \setcounter{CAPETannexe}{0}%
    \loop\ifnum\value{CAPETannexe}<\value{CAPETexo}%
      \stepcounter{CAPETannexe}%
      \edef\CAPETcurnum{\theCAPETannexe}%
      \ifcsname CAPETnotionstart\CAPETcurnum\endcsname
        \par\medskip
        {\normalsize\bfseries\color{colExo}\csuse{CAPETnotionstart\CAPETcurnum}}%
        \par\nobreak\smallskip
      \fi
      \begin{tcolorbox}[exoBox, breakable,
        title={Corrige de l'exercice \CAPETcurnum\
               \hfill\normalfont\footnotesize
               \hyperref[capet:eno:\CAPETcurnum]{enonce p.~\pageref{capet:eno:\CAPETcurnum}}}]%
      \label{capet:cor:\CAPETcurnum}%
      \csuse{CAPETcorrige\CAPETcurnum}%
      \end{tcolorbox}%
    \repeat
  }
\else
  % --- Modes historiques, comportement strictement inchange ---
  \newcommand{\CAPETSHOW}[2]{\ifdefined\CAPETCORRIGE #2\else #1\fi}
  \newcommand{\CAPETANNEXE}{}
\fi

% ── Difficulté 1 à 3 (étoiles pleines, amssymb) ───────────────────────
\newcommand{\difficulte}[1]{%
  \ifcase#1\relax
  \or $\bigstar$\;niveau 1/3%
  \or $\bigstar\bigstar$\;niveau 2/3%
  \or $\bigstar\bigstar\bigstar$\;niveau 3/3%
  \fi
}

% ── Style encadré extrait de document technique (banque d'exercices) ──
\tcbset{
  docTech/.style={
    breakable,
    enhanced,
    colback=colGrayLight,
    colframe=colSeparator,
    fonttitle=\bfseries\small\itshape,
    coltitle=black,
    colbacktitle=colGrayLight,
    left=4pt, right=4pt, top=3pt, bottom=3pt,
    before skip=4pt, after skip=4pt,
  }
}

% ── Style tcolorbox de section principale ────────────────────────────
% Refonte 2026 : bandeau titre coloré arrondi + fond très pâle, sans cadre.
\tcbset{
  ficheSection/.style={
    breakable, enhanced jigsaw,
    frame hidden, boxrule=0pt, arc=1.8mm,
    colback=#1!4!white,
    colbacktitle=#1,
    coltitle=white,
    fonttitle=\sffamily\bfseries,
    left=7pt, right=7pt, top=5pt, bottom=5pt,
    toptitle=3pt, bottomtitle=3pt,
    before skip=10pt, after skip=8pt,
  }
}

% ── Style encadré question type ──────────────────────────────────────
% Carte blanche, bandeau sombre, liseré très clair.
\tcbset{
  questiontype/.style={
    breakable, enhanced jigsaw,
    colback=white,
    frame hidden, boxrule=0pt, arc=1.8mm,
    colbacktitle=colElec!78!black,
    coltitle=white,
    fonttitle=\sffamily\bfseries\small,
    left=7pt, right=7pt, top=4pt, bottom=4pt,
    toptitle=2.5pt, bottomtitle=2.5pt,
    before skip=8pt, after skip=8pt,
    borderline={0.6pt}{0pt}{colElec!12!white},
  }
}

% ── Style encadré jury ───────────────────────────────────────────────
% Bandeau ambre, fond crème, sans cadre.
\tcbset{
  jury/.style={
    breakable, enhanced jigsaw,
    colback=colJuryBg,
    frame hidden, boxrule=0pt, arc=1.8mm,
    colbacktitle=colJuryFrame,
    coltitle=white,
    fonttitle=\sffamily\bfseries,
    left=7pt, right=7pt, top=4pt, bottom=4pt,
    toptitle=3pt, bottomtitle=3pt,
    before skip=10pt, after skip=8pt,
  }
}

% ── En-tête / pied de page communs (refonte 2026) ─────────────────────
% #1 code fiche, #2 titre court, #3 couleur domaine
\newcommand{\ficheHeaders}[3]{%
  \fancyhead[L]{\small\sffamily{\bfseries\textcolor{#3}{#1}}\hspace{0.7em}\textcolor{colSlate}{#2}}%
  \fancyhead[R]{\small\sffamily\textcolor{colSlate}{CAPET SII IE\, --- \,Session 2026}}%
  \fancyfoot[C]{\small\sffamily\textcolor{colSlate}{\thepage\,/\,\pageref{LastPage}}}%
  \fancyhead[C]{}\fancyfoot[L]{}\fancyfoot[R]{}%
}

% ── Bloc en-tête de fiche (BLOC 0, refonte 2026) ──────────────────────
% #1 code  #2 titre  #3 couleur domaine  #4 domaine  #5 ligne présence  #6 sous-thèmes
\newcommand{\ficheEntete}[6]{%
  \begin{tcolorbox}[enhanced, breakable,
    frame hidden, boxrule=0pt, arc=2.5mm,
    interior style={left color=#3, right color=#3!75!black},
    left=9pt, right=9pt, top=7pt, bottom=7pt,
    before skip=0pt, after skip=6pt]
    {\color{white}\sffamily\bfseries\LARGE #1\quad #2}\\[4pt]
    {\color{#3!22!white}\sffamily\small \MakeUppercase{#4}\quad\textbar\quad #5}\\[1pt]
    {\color{#3!35!white}\sffamily\footnotesize Sous-thèmes : #6}
  \end{tcolorbox}%
}

% ── Mise en page ─────────────────────────────────────────────────────
\pagestyle{fancy}
\fancyhf{}
\renewcommand{\headrulewidth}{0.5pt}
\renewcommand{\footrulewidth}{0pt}
\AtBeginDocument{%
  \renewcommand{\headrule}{{\color{colRule}\hrule width\headwidth height 0.5pt}}%
}

\setlist[itemize]{leftmargin=1.2em, itemsep=1pt, topsep=2pt, parsep=0pt}
\setlist[enumerate]{leftmargin=1.2em, itemsep=3pt, topsep=2pt, parsep=0pt}

\setlength{\parskip}{2pt}
\setlength{\columnsep}{10pt}

\def\CAPETBEGINFILE{\begin{document}}
\def\CAPETENDFILE{\end{document}}
\else
\def\CAPETBEGINFILE{}
\def\CAPETENDFILE{}
\fi

\ficheHeaders{T13}{Machine asynchrone (MAS)}{colElec}

\CAPETBEGINFILE

% -----------------------------------------------------------------------
%  BLOC EN-TETE
% -----------------------------------------------------------------------

\ficheEntete{T13}{Machine asynchrone (MAS)}{colElec}%
  {électrotechnique - Machines électriques}%
  {\textbf{Présence :} \textbf{7/13} sujets (2013-2026) - thème \emph{fréquent}}%
  {vitesse synchrone et glissement - bilan de puissance (Pa - Pt - Pméca - Pu) - couple électromagnétique - démarrage étoile-triangle - commande V/f (voir T14)}

% -----------------------------------------------------------------------
%  1. RAPPELS DE COURS
% -----------------------------------------------------------------------

\begin{tcolorbox}[
  ficheSection=colElec,
  title={1. Rappels de cours},
]

% -- 1.1 Vitesse synchrone et glissement ------------------------------
\textbf{1.1 Vitesse de synchronisme et glissement}

\smallskip
\begin{center}
\renewcommand{\arraystretch}{1.5}
\begin{tabular}{lll}
\toprule
\textbf{Grandeur} & \textbf{Expression} & \textbf{Unité} \\
\midrule
Vitesse synchrone & $N_s = \dfrac{60\,f}{p}$ & tr\,min$^{-1}$ \\[6pt]
Glissement & $g = \dfrac{N_s - N}{N_s} = \dfrac{\omega_s - \omega}{\omega_s}$ & - \\[6pt]
Vitesse pulsation rotorique & $\omega_r = g\,\omega_s$ & rad\,s$^{-1}$ \\
\bottomrule
\end{tabular}
\end{center}
\textit{Méthode : calculer $p = 60f/N_{nom}$, arrondir à l'entier le plus proche, puis recalculer $N_s$.}

\medskip\hrule\vspace{4pt}

% -- 1.2 Bilan de puissance --------------------------------------------
\textbf{1.2 Bilan de puissance - flux de droite à gauche}

\smallskip
\[
  P_a \xrightarrow{-P_{Js}-P_{Fs}} P_t \xrightarrow{-P_{Jr} = -g P_t} P_{meca} = (1-g)P_t \xrightarrow{-P_{pm}} P_u = C_u\Omega
\]

\begin{center}
\renewcommand{\arraystretch}{1.4}
\begin{tabular}{lll}
\toprule
\textbf{Puissance} & \textbf{Expression} & \textbf{Remarque} \\
\midrule
Absorbée au stator & $P_a = 3\,V_s I_s \cos\varphi$ & $V_s$ = tension \emph{simple} \\
Pertes Joule stator & $P_{Js} = 3\,R_s I_s^2$ & souvent négligées \\
Transmise au rotor & $P_t = P_a - P_{Js} - P_{Fs}$ & puissance de l'entrefer \\
Pertes Joule rotor & $P_{Jr} = g\,P_t = 3 R'_R I'^2_r$ & \\
Puissance mécanique & $P_{meca} = (1-g)\,P_t$ & développée au rotor \\
Puissance utile & $P_u = P_{meca} - P_{pm} = C_u\,\Omega$ & sortie d'arbre \\
\bottomrule
\end{tabular}
\end{center}

\medskip\hrule\vspace{4pt}

% -- 1.3 Couple electromagnetique -------------------------------------
\textbf{1.3 Couple électromagnétique}

\smallskip
\[
  C_{em} = \frac{P_t}{\Omega_s}
  \qquad;\qquad
  C_u = C_{em} - C_p
  \qquad;\qquad
  C_{em} \approx \frac{3p\,V_s^2}{R'_R\,\omega_s^2} \cdot g \quad \text{(glissement faible)}
\]
Forme linéarisée : $C_{em} = A \cdot \omega_r$ avec $A = \dfrac{3p V_s^2}{R'_R \omega_s^2}$ - valable pour $g \ll 1$.

\smallskip
\textbf{Courbe couple-vitesse (zone linéarisée au voisinage du synchronisme) :}
\begin{center}
\begin{tikzpicture}[scale=0.95, font=\small]
  \draw[->] (0,0) -- (5.6,0) node[right] {$\Omega$};
  \draw[->] (0,-0.3) -- (0,2.5) node[above] {$C_{em}$};
  \draw[dashed] (4,0) -- (4,2.2);
  \node[below] at (4,0) {$\Omega_s$};
  \draw[thick, colElec] (1.2,2.2) -- (4,0);
  \node[colElec, above left] at (1.5,2.1) {\scriptsize zone linéarisée ($g\ll1$)};
\end{tikzpicture}
\end{center}
Au voisinage du synchronisme ($\Omega \to \Omega_s$, $g \to 0$), le couple électromagnétique décroît linéairement avec la vitesse : c'est cette portion, seule, que valide l'approximation $C_{em} = A\cdot\omega_r$.

\medskip\hrule\vspace{4pt}

% -- 1.4 Modele de Laplace et demarrage -------------------------------
\textbf{1.4 Équation mécanique et démarrage étoile-triangle}

\smallskip
Équation mécanique dans le domaine de Laplace :
\[
  \Omega(p) = \frac{K_{meca}}{1 + \tau_{meca}\,p}\,\bigl[C_{em}(p) - C_r(p)\bigr]
  \qquad
  K_{meca} = \frac{1}{f}\,,\quad \tau_{meca} = \frac{J}{f}
\]

\textbf{Démarrage étoile-triangle :} en étoile, chaque enroulement reçoit $V_{reseau}/\sqrt{3}$ - courant de démarrage divisé par 3 - couple de démarrage divisé par 3. Passage en triangle à quasi-vitesse nominale.

\end{tcolorbox}

% -----------------------------------------------------------------------
%  2. FORMULAIRE SYNTHETIQUE
% -----------------------------------------------------------------------

\begin{tcolorbox}[
  ficheSection=colElec,
  title={2. Formulaire synthétique},
]
\small
\begin{multicols}{2}
\renewcommand{\arraystretch}{1.45}
\begin{tabular}{@{}p{4.5cm}ll@{}}
\toprule
\textbf{Formule} & \textbf{Usage} & \textbf{U.} \\
\midrule
\multicolumn{3}{@{}l}{\footnotesize\textit{Vitesse et glissement}} \\
$N_s = 60f/p$ & Synchro & tr/min \\
$g = (N_s-N)/N_s$ & Glissement & - \\
$\omega_r = g\,\omega_s$ & Puls. rotorique & rad/s \\
\midrule
\multicolumn{3}{@{}l}{\footnotesize\textit{Puissances}} \\
$P_a = 3\,V_s I_s\cos\varphi$ & Absorbée & W \\
$P_{Jr} = g\,P_t$ & Joule rotor & W \\
$P_{meca} = (1-g)\,P_t$ & Mécanique & W \\
$P_u = P_{meca} - P_{pm}$ & Utile & W \\
\bottomrule
\end{tabular}

\columnbreak

\renewcommand{\arraystretch}{1.45}
\begin{tabular}{@{}p{4.5cm}ll@{}}
\toprule
\textbf{Formule} & \textbf{Usage} & \textbf{U.} \\
\midrule
\multicolumn{3}{@{}l}{\footnotesize\textit{Couple}} \\
$C_{em} = P_t/\Omega_s$ & Couple em & N\,m \\
$C_u = C_{em} - C_p$ & Couple utile & N\,m \\
$C_{em} \approx A\cdot\omega_r$ & Linéarisé & N\,m \\
\midrule
\multicolumn{3}{@{}l}{\footnotesize\textit{Modèle Laplace}} \\
$K_{meca} = 1/f$ & Gain méc. & - \\
$\tau_{meca} = J/f$ & Constante & s \\
\midrule
\multicolumn{3}{@{}l}{\footnotesize\textit{Données 2024}} \\
$p = 2$, $N_s = 1500$ tr/min & Extracteur & - \\
$g = 6{,}67\%$, $P_u = 100$ kW & Extracteur & - \\
\bottomrule
\end{tabular}
\end{multicols}
\normalsize
\end{tcolorbox}

% -----------------------------------------------------------------------
%  3. QUESTIONS TYPES
% -----------------------------------------------------------------------

\begin{tcolorbox}[
  ficheSection=colElec,
  title={3. Questions types - extraites des sujets réels},
]

% -- QT1 -------------------------------------------------------------

\begin{tcolorbox}[questiontype,
  title={QT1 - Vitesse synchrone, glissement, puissance absorbée \hfill\normalfont\footnotesize E2 disc. 2024 (A1-A7, extracteur d'air)}
]
\textbf{Énoncé.}
Plaque signalétique du moto-extracteur : $V_{composee} = 400\,\mathrm{V}$, $V_s = 230\,\mathrm{V}$, $I_s = 229\,\mathrm{A}$, $\cos\varphi_{mot} = 0{,}85$, $f = 50\,\mathrm{Hz}$, $N = 1400\,\mathrm{tr/min}$, $p = 2$.
Calculer $N_s$, $g$ et $P_a$.

\tcbline

\textbf{Corrigé.}

\textbf{Étape 1 — Vitesse synchrone.}
\[
  N_s = \frac{60 \times 50}{2} = \mathbf{1500\,\mathrm{tr/min}}
\]

\textbf{Étape 2 — Glissement.}
\[
  g = \frac{N_s - N}{N_s} = \frac{1500 - 1400}{1500} = 0{,}0\overline{6} \approx \mathbf{6{,}67\,\%}
\]

\textbf{Étape 3 — Puissance absorbée.}
\[
  P_a = 3\,V_s I_s \cos\varphi = 3 \times 230 \times 229 \times 0{,}85 = 134\,308{,}5\,\si{\watt} \approx \mathbf{134{,}3\,\si{\kilo\watt}}
\]

\textbf{Conclusion.} Un glissement de $6{,}67\,\%$ est cohérent avec un fonctionnement en charge normale d'une MAS (typiquement 3 à 8\,\%) : la machine n'est ni à l'arrêt ni en survitesse anormale. \textit{Rappel : utiliser la tension \emph{simple} $V_s = 230\,\mathrm{V}$ (et non 400\,V) avec le facteur 3.}
\end{tcolorbox}

\vspace{4pt}

% -- QT2 -------------------------------------------------------------

\begin{tcolorbox}[questiontype,
  title={QT2 - Bilan de puissance complet \hfill\normalfont\footnotesize E2 disc. 2024 (A5-A7, extracteur)}
]
\textbf{Énoncé.}
Données DT1 : $R'_R = 0{,}1\,\Omega$, $I'_r = 122\,\mathrm{A}$, $g = 0{,}067$, $C_p = 200\,\mathrm{N\!\cdot\!m}$, $N = 1400\,\mathrm{tr/min}$.
Calculer $P_t$, $P_{Jr}$, $P_{meca}$ et $P_u$.

\tcbline

\textbf{Corrigé.}

\textbf{Étape 1 — Puissance transmise.}
\[
  P_t = \frac{3 R'_R}{g} I'^2_r = \frac{3 \times 0{,}1}{0{,}067} \times 122^2 = 66\,645\,\si{\watt} \approx 66{,}6\,\si{\kilo\watt}
\]

\textbf{Étape 2 — Pertes Joule rotor et puissance mécanique.}
\[
  P_{Jr} = g\,P_t = 0{,}067 \times 66{,}6 \approx 4{,}5\,\si{\kilo\watt}
  \qquad;\qquad
  P_{meca} = (1-g)\,P_t = 0{,}933 \times 66{,}6 \approx 62{,}2\,\si{\kilo\watt}
\]

\textbf{Étape 3 — Puissance utile.}
\[
  \Omega = \frac{1400 \times 2\pi}{60} = 146{,}6\,\si{\radian\per\second}
  \quad;\quad
  P_{pm} = C_p\,\Omega = 200 \times 146{,}6 = 29{,}3\,\si{\kilo\watt}
\]
\[
  P_u = P_{meca} - P_{pm} = 62{,}2 - 29{,}3 = \mathbf{32{,}9\,\si{\kilo\watt}}
\]

\textbf{Conclusion.} Chaque étage du bilan (Joule rotor, pertes mécaniques) prélève une part identifiable de $P_t$ : c'est cette décomposition, et non le seul résultat final $P_u$, que le jury attend pour juger la maîtrise du bilan de puissance. \textit{Note [corrigé reconstitué] : application numérique reconstituée à partir de données partielles du DT1 ; le sujet réel aboutit à $P_u \approx \SI{100}{\kilo\watt}$ avec l'ensemble des données du DT1. La méthode est identique.}
\end{tcolorbox}

\vspace{4pt}

% -- QT3 -------------------------------------------------------------

\begin{tcolorbox}[questiontype,
  title={QT3 - Couple EM linéarisé et modèle Laplace \hfill\normalfont\footnotesize E2 disc. 2024 (C1-C3)}
]
\textbf{Énoncé.}
Montrer que pour $g \approx 0$ : $C_{em} = \dfrac{3p\,V_s^2}{R'_R\,\omega_s^2}\cdot g$.
En déduire la fonction de transfert mécanique $\Omega(p)$ avec $J = 0{,}035\,\mathrm{kg\!\cdot\!m^2}$ et $f = 4 \times 10^{-3}\,\mathrm{N\!\cdot\!m\!\cdot\!s\!\cdot\!rad^{-1}}$.

\tcbline

\textbf{Corrigé.}

\textbf{Étape 1 — Approximation faible glissement.} Pour $g \ll 1$ : $I'_r \approx V_s g / R'_R$ et $P_t = 3(R'_R/g)I'^2_r = 3V_s^2 g/R'_R$. Avec $\Omega_s = \omega_s/p$ :
\[
  C_{em} = \frac{P_t}{\Omega_s} = \frac{p\,P_t}{\omega_s} = \frac{3p\,V_s^2 g}{R'_R \omega_s} = \underbrace{\frac{3p V_s^2}{R'_R \omega_s^2}}_{A} \cdot \underbrace{g\omega_s}_{\omega_r} = A\cdot\omega_r
\]

\textbf{Étape 2 — Modèle Laplace.} $J\,d\Omega/dt + f\,\Omega = C_{em} - C_r \;\Rightarrow\; \Omega(p) = \dfrac{K_{meca}}{1+\tau_{meca}\,p}\,[C_{em}(p)-C_r(p)]$
\[
  K_{meca} = \frac{1}{f} = \frac{1}{4\times10^{-3}} = 250\,\mathrm{rad\!\cdot\!s^{-1}\!\cdot\!N^{-1}\!\cdot\!m^{-1}}
  \qquad;\qquad
  \tau_{meca} = \frac{J}{f} = \frac{0{,}035}{0{,}004} = 8{,}75\,\mathrm{s}
\]

\textbf{Conclusion.} La forme linéarisée $C_{em} = A\cdot\omega_r$ transforme la MAS en un système du premier ordre standard, directement utilisable pour la synthèse d'un correcteur de vitesse — c'est l'intérêt principal de cette linéarisation autour du synchronisme.
\end{tcolorbox}

\end{tcolorbox}

% -----------------------------------------------------------------------
%  4. ATTENTES DU JURY
% -----------------------------------------------------------------------

\begin{tcolorbox}[
  jury,
  title={4. Attentes du jury - Erreurs récurrentes et points de vigilance},
]

\begin{enumerate}[label=\textbf{\arabic*.}, leftmargin=*, itemsep=5pt]

  \item \textbf{Facteur 3 en triphasé - erreur très fréquente.}
  $P_a = \mathbf{3}\,V_s I_s\cos\varphi$ et non $V_s I_s\cos\varphi$. Utiliser la tension \emph{simple} $V_s = 230\,\mathrm{V}$ (et non la composée 400\,V) si le facteur 3 est présent. Rapport jury 2024 : lacune signalée.

  \item \textbf{$p$ entier obligatoire.}
  Calculer $p = 60f/N_{nom}$, puis arrondir à l'entier le plus proche, puis recalculer $N_s$.

  \item \textbf{Glissement en décimal dans les formules.}
  $g = 0{,}067$ (décimal) dans $P_{Jr} = g\,P_t$, pas $g = 6{,}67$.

  \item \textbf{$P_t \neq P_{meca}$.}
  $P_t$ est la puissance transmise à l'entrefer (puissance de l'entrefer). $P_{meca} = (1-g)P_t$ est la puissance mécanique développée. $P_{Jr} = g\,P_t$ est perdue par effet Joule au rotor.

  \item \textbf{Démarrage étoile-triangle.}
  Le couple de démarrage est divisé par 3 en étoile - démarrage à charge réduite uniquement. Ne pas conclure que c'est adapté à toutes les charges.

  \item \textbf{Écrire l'expression littérale puis l'application numérique.}
  Rapport jury 2024 : "les résultats donnés directement, sans calcul et sans justification ne répondent pas aux attendus". Trois lignes par calcul minimum.

\end{enumerate}

\medskip
\begin{tcolorbox}[
  colback=colJuryFrame!12!colJury,
  colframe=colJuryFrame,
  left=4pt, right=4pt, top=2pt, bottom=2pt,
  fontupper=\small\itshape
]
Rapport jury 2024 : «~L'objectif de cette partie est de dimensionner le variateur et du moto-extracteur électrique pour optimiser la consommation énergétique de l'extracteur d'air. Seulement 55\,\% des candidats l'ont traitée correctement.~»\\
Rapport jury 2024 : «~Le jury regrette le manque de maîtrise des formules de base sur les réseaux triphasés, le régime sinusoïdal et les relations d'électrocinétique.~»
\end{tcolorbox}

\end{tcolorbox}

% -----------------------------------------------------------------------
%  PIED - PRESENCE DANS LES SUJETS
% -----------------------------------------------------------------------

\begin{tcolorbox}[
  ficheSection=colPresence,
  title={Présence dans les sujets analysés (2013-2026)},
]
\begin{center}
\footnotesize\renewcommand{\arraystretch}{1.15}
\begin{tabular}{lccccccccccccc}
\toprule
\textbf{Année} & 13 & 14 & 15 & 16 & 17 & 18 & 19 & 20 & 21 & 22 & 23 & 24 & 25 \\
\midrule
E1 (commune)   & - & - & - & - & - & - & - & $\checkmark$ & - & - & - & $\checkmark$ & - \\
E2 discipl. IE & - & - & - & $\checkmark$ & $\checkmark$ & $\checkmark$ & - & $\checkmark$ & $\checkmark$ & - & $\checkmark$ & $\checkmark$ & - \\
\bottomrule
\end{tabular}
\end{center}
\smallskip
\footnotesize\textbf{Sous-thèmes fréquents :}
bilan de puissance $P_a \to P_t \to P_{meca} \to P_u$ (7/13) - vitesse synchrone et glissement (7/13) - couple EM (5/13) - modèle Laplace mécanique (2/13) - démarrage étoile-triangle (3/13)

\medskip
\textbf{Ordres de grandeur (repères rapides).}
\begin{itemize}[leftmargin=1.3em]
  \item Réseau BT FR : tension phase-neutre typique \SI{230}{\volt}, tension phase-phase typique \SI{400}{\volt}, fréquence \SI{50}{\hertz}.
  \item Courants d'étude fréquents en sujet : de quelques \si{\ampere} à quelques dizaines de \si{\ampere}.
  \item En conversion d'énergie, un rendement attendu reste généralement entre 0{,}8 et 0{,}98 selon la technologie.
\end{itemize}

\medskip
\textbf{Synthèse express (à refaire sans support).}
\begin{enumerate}[leftmargin=1.4em]
  \item Poser les hypothèses du modèle et les conventions de signe.
  \item Écrire 2 à 4 relations clés avec unités.
  \item Réaliser un calcul numérique complet et vérifier l'ordre de grandeur.
  \item Justifier le résultat par un argument technique de niveau CAPET.
\end{enumerate}

\end{tcolorbox}

\CAPETENDFILE
